% Evaluation — v0.2 prose.
% Target: ~4 pages. The core empirical section. Every number cited
% sources from docs/evaluation.md; the per-detector table is
% inline rather than appendix-deferred.

\section{Evaluation}
\label{sec:evaluation}

\subsection{Research questions}

The corpus + measurement pipeline is built to answer four
questions about the detector suite, each falsifiable against the
labeled data:

\begin{description}
  \item[RQ1.] Per detector, what are precision, recall, and
    F\textsubscript{1}\,?
  \item[RQ2.] Do the detectors generalise beyond the samples
    used during detector iteration?
  \item[RQ3.] Do two independent implementations of the same
    detector produce verdicts that agree byte-for-byte?
  \item[RQ4.] Where does the measurement fall short?
\end{description}

Sections~\ref{subsec:eval-results} through
\ref{subsec:eval-limitations} report each in turn. The numbers
in this section source from
\verb|docs/evaluation.md|~\cite{primitivehub:2026:eval} and the
JSON baselines under \verb|bench/baseline/|; the discipline
documented in
\texttt{paper/README.md}~\cite{primitivehub:2026:paper-readme}
forbids invented numbers, so every quantitative claim here can
be traced back through the source-of-truth files to the
detector + corpus combination that generated it.

\subsection{Corpus}

% Source: docs/corpus-v0.1.md; bench/corpus/manifest.yaml.
The Adversarial Corpus v0.2 ships 109 labeled samples, split
into 99 in-sample (used during detector iteration) and 10
holdout (reserved as the overfit detector). Sample distribution
across labels: 52 adversarial in-sample plus 10 adversarial
holdout (62 total adversarial), 37 benign in-sample (no benign
holdouts in v0.2 by design --- the holdout's purpose is
overfit-recall not overfit-precision), and 10 borderline
in-sample (where reasonable labelers could mark either way).
Per-category support ranges from 30 to 45 samples in the
in-sample slice; the holdout slice carries 1 sample per
detector after Round~9's coverage closure.

Labeling discipline is documented in
\verb|docs/corpus-v0.1.md|~\cite{primitivehub:2026:corpus-spec}.
Each sample carries a primary CCI category, an expected verdict
($\in$ \{PASS, WARN, FAIL\}), and a textual provenance note
(\verb|incident_reference: synthesized| when not drawn from an
observed in-the-wild incident). Authoring and review are by repo
maintainers; no crowdsourced labels in v0.2. The corpus is
append-only; sample removal requires an
ADR-level justification.

The corpus license is CC-BY-4.0 (corpus content) so external
researchers can adapt the labeled samples in their own work and
cite this paper as the source.

\subsection{Methodology}

% Source: bench/scripts/measure.{ts,py}; docs/decisions/0032.
Per-sample, per-detector scoring follows the standard binary-
classification convention adapted to the multi-CCI setting:

\begin{itemize}
  \item For the sample's \emph{primary} detector (the detector
        whose \verb|cci_id| matches the sample's primary
        category):
    \begin{itemize}
      \item adversarial or borderline sample + verdict severity
            $\geq$ expected $\Rightarrow$ TP
      \item adversarial or borderline + verdict severity
            $<$ expected $\Rightarrow$ FN
      \item benign + verdict severity $\geq$ WARN
            $\Rightarrow$ FP
      \item benign + verdict severity = PASS
            $\Rightarrow$ TN
    \end{itemize}
  \item For \emph{other} detectors (cross-trip), only the
        FP/TN dimension applies (we have no ground truth for
        what detector X should do on detector Y's adversarial
        content):
    \begin{itemize}
      \item benign + flagged $\Rightarrow$ FP on the other
            detector
      \item benign + clean $\Rightarrow$ TN on the other
            detector
    \end{itemize}
\end{itemize}

Per-detector precision, recall, and F\textsubscript{1} use the
standard formulas. Macro-averaging is unweighted across the
ten detectors (no class-imbalance argument applies because each
CCI category is a separate defense, not a sample-frequency
class).

The holdout slice is measured and reported separately. Per
ADR-0038~\cite{primitivehub:2026:adr-0038}, holdout
F\textsubscript{1} is informational --- the regression gate does
not enforce holdout regressions because the gate is what would
let the detectors be tuned against the holdout (defeating its
purpose). Instead, the holdout closure round (Round 7 PR 7.1
plus Round 10 PR 10.1 / 10.2 / 10.3 in this paper's git history)
ratchets the detectors to close the holdout gaps as a separate
exercise, with each closure recorded in the FN ledger.

Cross-language parity: every measurement runs both Py and TS
engines against the same sample bytes. The parity check at
\verb|bench/scripts/parity_check.py| compares per-detector
\verb|(tp, fp, tn, fn)| tuples and exits non-zero on any drift.

\subsection{Results}
\label{subsec:eval-results}

% Source: docs/evaluation.md (auto-generated).
The headline result, against the v0.2 corpus as of submission:
macro F\textsubscript{1} = 1.000 on the 99-sample in-sample
slice and 1.000 on the 10-sample holdout slice, with precision
and recall individually 1.000 on every detector that has a
sample in the relevant slice, and cross-language verdict-count
parity intact for every detector on every sample.

Table~\ref{tab:per-detector-in-sample} reports the per-detector
breakdown on the in-sample slice. Both engines produce
byte-identical \verb|(tp, fp, tn, fn)| tuples, so we report a
single column.

\begin{table}[h]
\centering
\small
\begin{tabular}{lrrrrrrrr}
\toprule
Detector & Precision & Recall & F\textsubscript{1} & TP & FP & TN & FN & Support \\
\midrule
CCI-001 & 1.000 & 1.000 & 1.000 &  4 & 0 & 36 & 0 & 41 \\
CCI-002 & 1.000 & 1.000 & 1.000 &  3 & 0 & 37 & 0 & 40 \\
CCI-003 & 1.000 & 1.000 & 1.000 &  7 & 0 & 37 & 0 & 44 \\
CCI-004 & 1.000 & 1.000 & 1.000 &  7 & 0 & 37 & 0 & 44 \\
CCI-005 & 1.000 & 1.000 & 1.000 &  5 & 0 & 37 & 0 & 42 \\
CCI-006 & 1.000 & 1.000 & 1.000 &  7 & 0 & 37 & 0 & 44 \\
CCI-007 & 1.000 & 1.000 & 1.000 &  8 & 0 & 37 & 0 & 45 \\
CCI-008 & 1.000 & 1.000 & 1.000 &  7 & 0 & 37 & 0 & 44 \\
CCI-009 & 1.000 & 1.000 & 1.000 &  8 & 0 & 37 & 0 & 45 \\
CCI-010 & 1.000 & 1.000 & 1.000 &  6 & 0 & 37 & 0 & 43 \\
\midrule
\textbf{Macro} & \textbf{1.000} & \textbf{1.000} & \textbf{1.000} & --- & --- & --- & --- & --- \\
\bottomrule
\end{tabular}
\caption{Per-detector in-sample precision, recall, and
F\textsubscript{1} on the v0.2 corpus (99-sample slice). Both
Py and TS engines report identical tuples; cross-language
parity intact for all 10 detectors.}
\label{tab:per-detector-in-sample}
\end{table}

Table~\ref{tab:per-detector-holdout} reports the holdout slice
after Round~10's three closure PRs. The 10-sample slice has 1
sample per detector following Round~9's coverage expansion;
ten of ten detectors report F\textsubscript{1} = 1.000.

\begin{table}[h]
\centering
\small
\begin{tabular}{lrrrrrr}
\toprule
Detector & Precision & Recall & F\textsubscript{1} & TP & FN & Support \\
\midrule
CCI-001 & 1.000 & 1.000 & 1.000 & 1 & 0 & 1 \\
CCI-002 & 1.000 & 1.000 & 1.000 & 1 & 0 & 1 \\
CCI-003 & 1.000 & 1.000 & 1.000 & 1 & 0 & 1 \\
CCI-004 & 1.000 & 1.000 & 1.000 & 1 & 0 & 1 \\
CCI-005 & 1.000 & 1.000 & 1.000 & 1 & 0 & 1 \\
CCI-006 & 1.000 & 1.000 & 1.000 & 1 & 0 & 1 \\
CCI-007 & 1.000 & 1.000 & 1.000 & 1 & 0 & 1 \\
CCI-008 & 1.000 & 1.000 & 1.000 & 1 & 0 & 1 \\
CCI-009 & 1.000 & 1.000 & 1.000 & 1 & 0 & 1 \\
CCI-010 & 1.000 & 1.000 & 1.000 & 1 & 0 & 1 \\
\midrule
\textbf{Macro} & \textbf{1.000} & \textbf{1.000} & \textbf{1.000} & --- & --- & --- \\
\bottomrule
\end{tabular}
\caption{Per-detector holdout precision, recall, and
F\textsubscript{1} on the v0.2 corpus (10-sample slice).
The slice carries 10 adversarial samples (1 per detector); no
benign samples in the holdout by design. Cross-language parity
intact.}
\label{tab:per-detector-holdout}
\end{table}

The detector versions referenced in the tables are:
CCI-001 v\tonever{}, CCI-002 v\ttwover{}, CCI-003 v\tthreever{},
CCI-004 v\tfourver{}, CCI-005 v\tfivever{},
CCI-006 v\tsixver{}, CCI-007 v\tsevenver{},
CCI-008 v\teightver{}, CCI-009 v\tninever{},
CCI-010 v\ttenver{}.

\subsection{What the holdout proves}

% Source: FN ledger entries FN-007 through FN-017.
The six holdout samples in the original v0.2 split (the four
zero-holdout detectors received their first holdout in Round~9
PR~9.3) were authored simultaneously with the FN ledger entries
that named the resolution vector each detector would need to
catch them. That ordering matters: the prediction was committed
before the closure was measured. Round~7 PR~7.1 closed
FN-007 (URL-encoded decode-and-run), FN-008 (WebSocket load-time
fetch), and FN-012 (US Patent fake-citation) by shipping the
predicted detector ratchets. Round~10 PR~10.1, 10.2, and 10.3
closed FN-014 (Homebrew install + Gatekeeper bypass), FN-015
(Cyrillic confusable on unsupported vendor), FN-010 (env-var
$\to$ query-string exfil), and FN-017 (markdown-table-density
bury). Each closure shipped the resolution vector named in its
ledger entry, with the F\textsubscript{1} delta recorded in the
closure log.

One holdout sample, S-108 (CCI-009 viral CC-BY-NC-SA license
clause), was caught by the v0.2.3 detector unaided. The
prediction was that T9 would miss the noncommercial-ShareAlike
variant; the measurement showed the detector's CC-BY-* pattern
generalised across the noncommercial sub-variant. We log this
as FN-016 (informational; no closure required) and treat it as
the strongest evidence we have that the detector generalises
beyond the in-sample tuning pool. Three of four PR~9.3
predictions correct and one detector-generalises-better-than-
predicted is exactly the balance the holdout discipline was
designed to surface.

The combined evidence --- in-sample F\textsubscript{1} = 1.000,
holdout F\textsubscript{1} = 1.000 after closures, the FN ledger
recording every prediction-and-closure pair --- is the empirical
basis for the claim that the in-sample numbers are not overfit
to the in-sample phrasing pool.

\subsection{Limitations}
\label{subsec:eval-limitations}

% Source: docs/evaluation.md § Limitations (verbatim).
The numbers in this paper are honest about what they do and do
not prove. We enumerate eight specific limitations:

\paragraph{1. Corpus size.} v0.2 ships 109 labeled samples with
per-detector support 30 to 45. F\textsubscript{1} = 1.000 on
that support is the best demonstration we can make at this
corpus size; it is not the same claim as F\textsubscript{1} =
1.000 at 30{,}000 samples. Corpus expansion to per-detector
support past 100 samples is the immediate next-round work.

\paragraph{2. English-language bias.} Adversarial samples are
authored primarily in English. Attacker phrasing in Russian,
Chinese, Arabic, and Spanish is uncovered. Multilingual corpus
expansion is deferred to v0.3.

\paragraph{3. Static detection only.} The corpus measures
detectors operating on primitive content (markdown files). It
does \emph{not} measure runtime-injection defense (the agent
fetching a malicious resource after signing) or training-time
attacks against the underlying LLM~\cite{Carlini:2024:webcrawl-poisoning}.
Each maps to a different defense layer in the Trust Runtime
stack; see \texttt{docs/owasp-llm-top10-crosswalk.md} for the
scope boundary.

\paragraph{4. Vendor list scope.} Generic Unicode confusables
(T7 v0.3.1) defend the 43 single-word vendor names in
\verb|_VENDORS|. The full UTS \#39 confusables.txt distribution
($\sim$6{,}000 entries) plus UAX \#31 identifier-domain rules
is the natural extension when the vendor list outgrows manual
curation.

\paragraph{5. Synthesized provenance.} Most adversarial
samples carry \verb|incident_reference: synthesized| --- composed
to exercise a specific attack shape rather than drawn from
observed in-the-wild incidents. Field-incident provenance grows
as the \verb|/check| public route accumulates field reports;
see Section~\ref{sec:related} for the operational mechanism.

\paragraph{6. Holdout coverage per detector.} After Round~9
PR~9.3 + Round~10 closures, every detector has exactly 1
holdout sample. Multiple holdouts per detector (matching the
30--45 in-sample density) is desirable; v0.3 corpus work targets
3 holdouts per detector.

\paragraph{7. Replay versus gate.}
\verb|bench/cli/replay.py| reports per-sample verdicts and is
intended for paper-reviewer reproducibility;
\verb|bench/scripts/gate.sh| runs the aggregate measurement +
parity check that the CI gate enforces. Both consume the same
corpus + detectors; the gate is what blocks merges, the replay
is what answers ``show me this specific finding.''

\paragraph{8. No third-party verification.} These numbers are
produced by the same team that wrote the detectors. The
corpus and gate are open-source and reproducible by design;
independent replication is invited via the
responsible-disclosure surface documented in
\verb|SECURITY.md|~\cite{primitivehub:2026:security}, and
ledger closure entries credit external reporters.

\subsection{Threats to validity}

\paragraph{Internal validity.} The same team authored both the
detectors and the labels. The mitigation is the append-only FN
ledger so the chain from prediction to closure is auditable; an
external reviewer can re-run each closure's measurement and
verify the F\textsubscript{1} delta from the entry. The corpus is
open-source; \verb|SECURITY.md| documents the independent-
replication path.

\paragraph{External validity.} The corpus is synthesized for
most adversarial samples; real-world primitive content will
surface attack shapes not in the v0.2 set. The mitigation is
the regression gate and FP/FN ledger discipline are designed
to absorb new shapes as they appear without re-architecting; a
field-reported FN becomes a corpus sample, becomes an FN
ledger entry, becomes a detector ratchet via a corpus-
measured closure round.

\paragraph{Construct validity.} F\textsubscript{1} measures
verdict-versus-manifest agreement. The manifest encodes
maintainers' beliefs about what should fire. The mitigation
is the borderline label (sample where reasonable labelers
disagree) plus the four-path FP/FN triage discipline that
forces every diagnosed finding to classify itself into one
of \{edit content, allow with ledger record, ratchet detector,
escalate to security review\}. Disagreement at the corpus
labeling layer is recoverable via Path B (sample relabel);
disagreement at the detector layer is recoverable via Path C
(detector ratchet); neither path is silent.
